% ============================================================================
% Supplementary Information SI-3: Reaction-Network Kinetics Solver
%
% Drop-in subsection fragment. Parent SI must provide: amsmath/amssymb,
% booktabs, siunitx (with \angstrom), and the tcolorbox environments
% {evaluation}, {concern}, {threshbox}. Bibliography handled at the
% parent level. (Parameter reference lives in si_parameters.tex.)
% ============================================================================

\subsection{Reaction-Network Kinetics Solver}
\label{si:kinetics}

The kinetics solver consumes the current reaction graph and integrates a mass-action ODE to a long-time steady state. The resulting steady-state composition is the distribution used by the generative-loop sampler (SI-\ref{si:generative}).

%-----------------------------------------------------------------------------
\subsubsection{Mass-action ODE}
\label{sec:kinetics-ode}
%-----------------------------------------------------------------------------

For $n_\text{species}$ species and $n_\text{reactions}$ reactions, the ODE is the standard mass-action system
\begin{equation}
    \frac{d\mathbf{y}}{dt} = \mathbf{S}_\text{net}\,\mathbf{r}(\mathbf{y}), \qquad r_j(\mathbf{y}) = k_j^+ \prod_i y_i^{\nu_{ij}^-} - k_j^- \prod_i y_i^{\nu_{ij}^+},
\end{equation}
where $\mathbf{S}_\text{net} = \mathbf{S}^+ - \mathbf{S}^-$ is the net stoichiometry matrix and $\boldsymbol{\nu}^\pm$ are the reactant/product stoichiometry exponents. The right-hand side and its analytical sparse Jacobian
\begin{equation}
    J_{ik} = \frac{\partial}{\partial y_k}\!\left[\sum_j S_{ij}\, r_j(\mathbf{y})\right]
\end{equation}
are implemented as Numba-JIT kernels (\texttt{numba\_kernels.py}, \texttt{@njit(cache=True)}) to keep the inner solver loop in compiled code. Jacobian sparsity follows from the stoichiometry pattern: $J_{ik} \neq 0$ only where species $k$ appears as a reactant or product of a reaction that touches species $i$.

%-----------------------------------------------------------------------------
\subsubsection{Rate constants}
\label{sec:kinetics-rates}
%-----------------------------------------------------------------------------

\paragraph{Eyring transition-state theory.}
For each discovered reaction the forward and backward rate constants are
\begin{equation}
    k = \min\!\Bigl(\tfrac{k_B T}{h}\,\exp(-E_a / k_B T),\ k_\text{diff}\Bigr), \qquad k_\text{diff} = \SI{e10}{M^{-1}s^{-1}},
\end{equation}
with $k_B = \SI{8.617e-5}{eV/K}$, $h = \SI{4.136e-15}{eV\,s}$, and the input barrier capped at $E_a \leq \SI{10}{eV}$ before the exponential (purely to avoid floating-point overflow; barriers above this are off-policy anyway). The diffusion-limit cap prevents the unit-prefactor Eyring rate from exceeding the physical bimolecular collision rate at the default exploration temperature of \SI{500}{K} --- without it, near-zero barriers would produce non-physical rates that dominate the Jacobian conditioning.

\paragraph{Barrier source priority.}
The barrier that feeds Eyring is chosen by a strict priority:
\begin{enumerate}
    \item PBE0/def2-TZVPP separated barriers (\texttt{barrier\_\{forward,backward\}\_separated\_pbe0}) if both forward and backward are available --- this is the DFT-refined ground truth from the CPU pipeline (\S\ref{sec:kinetics-dft}).
    \item Force-integrated in-box ML barriers (\texttt{barrier\_\{forward,backward\}}) along the IRC trajectory.
    \item Otherwise: the reaction is dropped entirely from the ODE.
\end{enumerate}
\emph{ML separated-energy barriers (computed as the difference of energy-head evaluations on two unrelated geometries) are never used} --- they are the least reliable variant: differencing two energy-head evaluations on unrelated geometries is noisier than either the DFT separated estimate or the ML force-integrated estimate.

\paragraph{Manual buffer equilibria.}
Acid--base equilibria (water autoionisation, CO$_2$ hydration, H$_2$CO$_3$ dissociation, proton solvation; SI-\ref{si:seed}) bypass Eyring entirely. Their rate constants are stored as literal values in the \texttt{Reaction} row (\texttt{manual\_k\_fwd}, \texttt{manual\_k\_bwd}) and applied directly at any temperature. These reactions are tagged \texttt{discovery\_method = "manual\_equilibrium"} and are layered into the ODE on top of the discovered reactions.

\paragraph{Filters applied at build time.}
Reactions are dropped from the ODE if (i)~either side contains dot-disconnected SMILES (van der Waals dimers stabilised by dispersion; non-chemical), (ii)~either barrier is below $-\SI{0.05}{eV}$ (negative-$E_a$ artefacts from PES/IRC landing on dianion or zwitterion minima below the separated reactant pair), or (iii)~both barriers exceed the cutoff (default \SI{10}{eV} --- effectively unreachable, just bloats the system). Reactions are also direction-agnostic deduplicated: $A + B \to C$ and $C \to A + B$ collapse into one bidirectional reaction keyed on the unordered pair of reactant/product multisets, taking the lowest $\min(E_a^\text{fwd}, E_a^\text{bwd})$. Storing both as separate mass-action reactions would double-count the net flux for the same elementary step.

The Eyring prefactor at $T = \SI{500}{K}$ gives $k_B T / h \approx \SI{1.04e13}{s^{-1}}$, so a $\SI{1.0}{eV}$ barrier corresponds to $\sim \SI{9e2}{s^{-1}}$ and a $\SI{2.0}{eV}$ barrier to $\sim \SI{8e-8}{s^{-1}}$. The \SI{500}{K} default is a sampling choice: at the higher temperature, more reactions are kinetically accessible within the integration window, producing a richer steady-state distribution for the generative sampler to draw from. Physical-temperature analyses use the same graph but a separate solve at the target temperature.


%-----------------------------------------------------------------------------
\subsubsection{DFT energy refinement}
\label{sec:kinetics-dft}
%-----------------------------------------------------------------------------

The sole purpose of the DFT pipeline in this system is to upgrade the barriers fed to the Eyring rates above. There is no DFT geometry optimisation, no DFT vibrational analysis, and no DFT-based dataset assembly: only single-point energies on the geometries already produced by the ML pipeline, used to derive PBE0 barriers that replace the ML in-box estimates whenever available.

\paragraph{Method.}
PBE0/def2-TZVPP single-points via PySCF (\texttt{dft.RKS} for closed shell, \texttt{dft.UKS} for open shell; spin guessed from electron count modulo 2). SCF convergence \texttt{conv\_tol} $= 10^{-9}$, \texttt{max\_cycle} $= 300$, with a Newton retry on non-convergence. All energies are returned in eV.

\paragraph{Per-reaction work.}
For each non-manual reaction the pipeline computes three single-point energies on the in-box geometry (same atom set throughout, since the IRC is performed in the merged-system box):
\begin{itemize}
    \item $E_\text{TS}^\text{DFT}$ at the TS conformer geometry,
    \item $E_R^\text{DFT}$ at the relaxed reactant frame (the final position of the IRC reactant-side trajectory),
    \item $E_P^\text{DFT}$ at the relaxed product frame (final position of the IRC product-side trajectory).
\end{itemize}
The in-box DFT barriers are then $\Delta E_\text{fwd}^\text{in-box} = E_\text{TS} - E_R$ and $\Delta E_\text{bwd}^\text{in-box} = E_\text{TS} - E_P$.

\paragraph{Per-compound reference energies.}
For each compound participating in any DFT-refined reaction, a single-point at the compound's \emph{lowest-energy minimum} (the PES-graph entry with the smallest ML energy) is computed and cached on the \texttt{Compound} row. The cache is keyed by \texttt{(method, basis)}; subsequent reactions touching the same compound reuse the value rather than recomputing. The \emph{separated} barriers --- the ones the kinetics solver actually consumes --- are derived as
\begin{equation}
    \Delta E_\text{fwd}^\text{sep} = E_\text{TS}^\text{DFT} - \!\!\sum_{c \in \text{reactants}}\!\! E_c^\text{DFT}, \qquad
    \Delta E_\text{bwd}^\text{sep} = E_\text{TS}^\text{DFT} - \!\!\sum_{c \in \text{products}}\!\! E_c^\text{DFT},
\end{equation}
where the sums respect reactant/product stoichiometry. Because the reference energies use the compound's lowest minimum (not the IRC-relaxed in-box geometry), the separated barriers absorb the inter-fragment binding-energy correction that the in-box variant misses.

\paragraph{No Hessians, no frequency check.}
The DFT pipeline does not compute Hessians or check TS character at DFT level. The reasoning is empirical: hybrid-functional Hessians fail on the same floppy-mode pathologies as ML Hessians (small spurious negative eigenvalues from low-frequency modes, near-degenerate hindered rotors), at $10^2$--$10^3\times$ the cost. Energies are where DFT provides leverage over ML; for spectrum verification, a dedicated tighter ML model or a CC-level audit on a small subset would be more cost-effective than per-reaction DFT Hessians.

\paragraph{Cost.}
A PBE0/def2-TZVPP single-point scales as $\mathcal{O}(N^{2}$--$N^{3})$ with system size; wall times in the production deployment range from $\sim$\SI{10}{s} for diatomic fragments to $\sim$\SI{5}{min} per geometry for $N \sim 13$--$20$ heavy atoms with the def2-TZVPP basis. Each reaction therefore costs three TS-sized single-points plus (amortised over the compound population) a handful of reference single-points. The pipeline runs on a separate pool of CPU workers (one work item $=$ one reaction) that drain a Postgres-backed queue independently of the GPU exploration workers, so DFT throughput does not bottleneck reaction discovery.

Refining only energies, on only the geometries the ML model already produced, is a deliberate cost--benefit trade. Re-optimising at DFT level would correct ML geometry errors but add $10$--$50\times$ the wall time per reaction and would require a DFT-level IRC re-run to keep the reactant/product matchings consistent. Empirically the ML geometry errors are small for the chemistry covered here (RMSD $\lesssim \SI{0.05}{\angstrom}$ to PBE0 minima for relaxed carbohydrate-space compounds), so the single-point-only protocol captures the bulk of the DFT correction at a small fraction of the cost. The compound-reference cache is load-bearing: without it, every reaction would re-compute energies for the same recurring small compounds (water, formaldehyde, etc.).


%-----------------------------------------------------------------------------
\subsubsection{Species universe and initial conditions}
\label{sec:kinetics-ic}
%-----------------------------------------------------------------------------

The species set is the union of (i)~every compound that appears as a reactant or product in any selected reaction, and (ii)~the seed species with prescribed initial concentrations ($\mathrm{H_2O}$, $\mathrm{CH_2O}$, $\mathrm{CO_2}$, $\mathrm{H_2}$, $\mathrm{H^+}$, $\mathrm{OH^-}$, $\mathrm{H_2CO_3}$, $\mathrm{HCO_3^-}$, $\mathrm{H_3O^+}$). Seeds are included even if no reaction touches them, so the buffer equilibria can self-consistently relax to their target pH without depending on whether a discovered reaction happens to mention water.

Every species in the system receives a uniform baseline of $c_0 = \SI{e-3}{M}$ at $t = 0$, with the seed species overriding the baseline (all at $\SI{e-3}{M}$ in the production deployment). The uniform baseline is essential: a runtime-discovered species starting at exactly zero would never react, and the steady-state distribution would be dominated entirely by the seeds.

A species-concentration noise floor of \SI{e-20}{M} is applied throughout post-processing: anything below this is treated as ``not present'' for entropy and distribution computations.

%-----------------------------------------------------------------------------
\subsubsection{Solver}
\label{sec:kinetics-solver}
%-----------------------------------------------------------------------------

The ODE is integrated to $t_\text{max} = \SI{e8}{s}$ ($> 3$ years --- well past any realistic reactor residence time) using PETSc's BDF time-stepper with Numba-JIT RHS and Jacobian kernels.

\paragraph{PETSc configuration.}
\begin{itemize}
    \item Time stepper: \texttt{TS} type \texttt{bdf}, exact-final-time interpolation, max $5\times 10^7$ steps.
    \item Initial step: $\Delta t_0 = \SI{e-10}{s}$. The chemistry's fastest timescale at \SI{500}{K} is set by the diffusion limit, $\sim k_\text{diff}^{-1} = \SI{e-10}{s}$ for bimolecular rates, so the initial step matches.
    \item Tolerances: $\texttt{atol} = \SI{e-16}{}$, $\texttt{rtol} = \SI{e-12}{}$. Very tight because some species sit at concentrations far below the noise floor and the solver's atol must clear them.
    \item Newton inner solver: \texttt{SNES} with a \texttt{KSP} preconditioner-only LU factorisation (\texttt{type=preonly}, \texttt{pc=lu}). The Jacobian factorisation is the expensive step; MUMPS is the default direct solver when available.
\end{itemize}

\paragraph{Output grid.}
The solver reports the full trajectory at $400$ geometrically spaced time points covering $[\SI{e-12}{s}, \SI{e8}{s}]$ for the frontend concentration plot, plus the canonical \emph{decade grid} of $19$ points at $10^{-10}, 10^{-9}, \ldots, 10^{8}$~s used by the kinetics snapshot. The decade grid spans every kinetically relevant order of magnitude, from femtosecond chemistry to multi-year residence times.
